% Options for packages loaded elsewhere
\PassOptionsToPackage{unicode}{hyperref}
\PassOptionsToPackage{hyphens}{url}
\documentclass{amsart}
\usepackage{dzg-unified}
% dzg-unified provides: amsmath, amsthm, amssymb, amsfonts, mathtools,
% mathrsfs, stmaryrd, bm, esint, upgreek, graphicx, xcolor[dvipsnames],
% tikz, tikz-cd, etoolbox, xparse, pgfornament, marginnote

\newif\ifshort
\shortfalse

\usepackage{hyperref}
\hypersetup{colorlinks,citecolor=blue,linktocpage,hyperindex=true,backref=true}
\newcommand\todoinline[1]{\todo[inline]{#1}}

\title{Research Draft}
\date{\today}


\usepackage[utf8]{inputenc}
\usepackage[english]{babel}
\usepackage{fontenc}

\newdimen\figrasterwd
\figrasterwd\textwidth
\floatplacement{figure}{H}

\newcolumntype{P}[1]{>{\centering\arraybackslash}p{#1}}

\makeatletter
\AddToHook{cmd/section/before}{\FloatBarrier}
\AddToHook{cmd/subsection/before}{\FloatBarrier}
\makeatother

\usetikzlibrary{arrows.meta, cd, fadings, patterns, calc, matrix, positioning, decorations, decorations.pathreplacing, decorations.markings, shapes, backgrounds, fit, shapes.geometric, intersections, hobby, matrix,arrows,decorations.pathmorphing}
\tikzfading[name=fade out, inner color=transparent!0, outer color=transparent!100]

% biblatex is provided by dzg-unified
\addbibresource{global.bib}

\renewcommand{\hat}{\widehat}
\renewcommand{\tilde}{\widetilde}

% Pandoc compatibility macros
\providecommand{\tightlist}{%
  \setlength{\itemsep}{0pt}\setlength{\parskip}{0pt}}

% Standard Pandoc macros for figures

\makeatletter
\providecommand{\pandocbounded}[1]{#1}
\makeatother



\begin{document}


\maketitle
\tableofcontents

\newpage

\section{Loops and suspension}\label{sec-loops-suspension}

Let \(\mathcal C\) be a pointed \(\infty\)-category, with zero object
\(*\), and suppose that it has the finite limits and colimits required
below.

\subsection{Fibers and cofibers}\label{sec-cobase-changes}

For \(f\colon X\to Y\), its fiber over the canonical basepoint
\(0\colon *\to Y\) is defined by the cartesian square

\begin{Shaded}
\begin{Highlighting}[]
\CommentTok{\%\%| filename: fiber{-}square}
\CommentTok{\%\%| additionalPackages: \textbackslash{}usepackage\{amsmath,amssymb,tikz{-}cd\}}
\KeywordTok{\textbackslash{}begin}\NormalTok{\{}\ExtensionTok{tikzcd}\NormalTok{\}}
\FunctionTok{\textbackslash{}operatorname}\NormalTok{\{Fib\}(f)}
  \FunctionTok{\textbackslash{}arrow}\NormalTok{[r,"p\_X"]}
  \FunctionTok{\textbackslash{}arrow}\NormalTok{[d,"p\_*"\textquotesingle{}]}
  \FunctionTok{\textbackslash{}arrow}\NormalTok{[dr,phantom,very near start,"}\FunctionTok{\textbackslash{}lrcorner}\NormalTok{"] \&}
\NormalTok{X }\FunctionTok{\textbackslash{}arrow}\NormalTok{[d,"f"]}\FunctionTok{\textbackslash{}\textbackslash{}}
\NormalTok{* }\FunctionTok{\textbackslash{}arrow}\NormalTok{[r,"0"\textquotesingle{}] \& Y}
\KeywordTok{\textbackslash{}end}\NormalTok{\{}\ExtensionTok{tikzcd}\NormalTok{\}}
\end{Highlighting}
\end{Shaded}

The fiber-product notation for its apex is \[
\operatorname{Fib}(f)=X\times_Y *.
\] The cofiber is defined by the cocartesian square

\begin{Shaded}
\begin{Highlighting}[]
\CommentTok{\%\%| filename: cofiber{-}square}
\CommentTok{\%\%| additionalPackages: \textbackslash{}usepackage\{amsmath,amssymb,tikz{-}cd\}}
\KeywordTok{\textbackslash{}begin}\NormalTok{\{}\ExtensionTok{tikzcd}\NormalTok{\}}
\NormalTok{X}
  \FunctionTok{\textbackslash{}arrow}\NormalTok{[r,"f"]}
  \FunctionTok{\textbackslash{}arrow}\NormalTok{[d,"0"\textquotesingle{}]}
  \FunctionTok{\textbackslash{}arrow}\NormalTok{[dr,phantom,very near end,"}\FunctionTok{\textbackslash{}ulcorner}\NormalTok{"] \&}
\NormalTok{Y }\FunctionTok{\textbackslash{}arrow}\NormalTok{[d,"i\_Y"]}\FunctionTok{\textbackslash{}\textbackslash{}}
\NormalTok{* }\FunctionTok{\textbackslash{}arrow}\NormalTok{[r,"i\_*"\textquotesingle{}] \& }\FunctionTok{\textbackslash{}operatorname}\NormalTok{\{Cof\}(f)}
\KeywordTok{\textbackslash{}end}\NormalTok{\{}\ExtensionTok{tikzcd}\NormalTok{\}}
\end{Highlighting}
\end{Shaded}

The pushout notation for its apex is \[
\operatorname{Cof}(f)=Y\amalg_X *.
\] These definitions specialize in pointed spaces to the usual homotopy
fiber and homotopy cofiber.

\subsection{Loops and suspension}\label{loops-and-suspension}

For an object \(X\) of \(\mathcal C\), loops are defined by the pullback
square

\begin{Shaded}
\begin{Highlighting}[]
\CommentTok{\%\%| filename: loop{-}square}
\CommentTok{\%\%| additionalPackages: \textbackslash{}usepackage\{amsmath,amssymb,tikz{-}cd\}}
\KeywordTok{\textbackslash{}begin}\NormalTok{\{}\ExtensionTok{tikzcd}\NormalTok{\}}
\FunctionTok{\textbackslash{}Omega}\NormalTok{ X}
  \FunctionTok{\textbackslash{}arrow}\NormalTok{[r,"}\FunctionTok{\textbackslash{}operatorname}\NormalTok{\{pr\}\_2"]}
  \FunctionTok{\textbackslash{}arrow}\NormalTok{[d,"}\FunctionTok{\textbackslash{}operatorname}\NormalTok{\{pr\}\_1"\textquotesingle{}]}
  \FunctionTok{\textbackslash{}arrow}\NormalTok{[dr,phantom,very near start,"}\FunctionTok{\textbackslash{}lrcorner}\NormalTok{"] \&}
\NormalTok{* }\FunctionTok{\textbackslash{}arrow}\NormalTok{[d,"0"]}\FunctionTok{\textbackslash{}\textbackslash{}}
\NormalTok{* }\FunctionTok{\textbackslash{}arrow}\NormalTok{[r,"0"\textquotesingle{}] \& X}
\KeywordTok{\textbackslash{}end}\NormalTok{\{}\ExtensionTok{tikzcd}\NormalTok{\}}
\end{Highlighting}
\end{Shaded}

Suspension is defined by the pushout square

\begin{Shaded}
\begin{Highlighting}[]
\CommentTok{\%\%| filename: suspension{-}square}
\CommentTok{\%\%| additionalPackages: \textbackslash{}usepackage\{amsmath,amssymb,tikz{-}cd\}}
\KeywordTok{\textbackslash{}begin}\NormalTok{\{}\ExtensionTok{tikzcd}\NormalTok{\}}
\NormalTok{X}
  \FunctionTok{\textbackslash{}arrow}\NormalTok{[r,"0"]}
  \FunctionTok{\textbackslash{}arrow}\NormalTok{[d,"0"\textquotesingle{}]}
  \FunctionTok{\textbackslash{}arrow}\NormalTok{[dr,phantom,very near end,"}\FunctionTok{\textbackslash{}ulcorner}\NormalTok{"] \&}
\NormalTok{* }\FunctionTok{\textbackslash{}arrow}\NormalTok{[d,"j\_1"]}\FunctionTok{\textbackslash{}\textbackslash{}}
\NormalTok{* }\FunctionTok{\textbackslash{}arrow}\NormalTok{[r,"j\_2"\textquotesingle{}] \& }\FunctionTok{\textbackslash{}Sigma}\NormalTok{ X}
\KeywordTok{\textbackslash{}end}\NormalTok{\{}\ExtensionTok{tikzcd}\NormalTok{\}}
\end{Highlighting}
\end{Shaded}

Thus \[
\Omega X=*\times_X *,
\qquad
\Sigma X=*\amalg_X *.
\] In pointed spaces these are the usual loop-space and
reduced-suspension constructions.

\subsection{Arrow categories and path
spaces}\label{sec-interval-presentations}

For a higher category \(C\), evaluation at the endpoints of the
constructed walking arrow gives \[
[[1],C]\longrightarrow C\times C.
\]

The fiber over \((x,y)\) is the hom-category \([x,y]_C\); applying
\(\Pi_\infty\) gives \(\operatorname{Map}_C(x,y)\). The fiber over
\((x,x)\) is the endomorphism category of \(x\), and its full
subcategory on the invertible objects is the automorphism category of
\(x\).

\subsection{The suspension-loop adjunction}\label{sec-adjunctions}

When \(\mathcal C\) has finite limits and colimits, suspension is left
adjoint to loops: \[
\adj{\mathcal C}{\mathcal C}{\Sigma}{\Omega},
\qquad \Sigma\dashv\Omega.
\] For \(\mathcal C=\mathcal S_*\), the \(\infty\)-category of pointed
spaces, this recovers the classical adjunction \[
\operatorname{Map}_*(\Sigma X,Y)\simeq
\operatorname{Map}_*(X,\Omega Y).
\]

\subsection{Fiber sequences}\label{fiber-sequences}

A composable pair \(F\to E\to B\) is a fiber sequence when \(F\) is
equivalent to the homotopy fiber over a specified basepoint of \(B\).
Applying homotopy groups gives the long exact sequence. Its
component-level portion is the pointed-set sequence recorded in
\textcite{sec-pi0-fiber}.

\clearpage
 \printbibliography

\end{document}
